\chapter*{Résumé}
\addcontentsline{toc}{chapter}{Résumé}

\textbf{Contexte.} L'accélération de la transformation numérique au Maroc a créé une demande d'intermédiation de confiance entre les consommateurs et les entreprises locales. Les places de marché multi-acteurs peuvent favoriser la participation économique, mais elles introduisent des obligations complexes en matière de gouvernance de la sécurité, de protection des données et de responsabilité des plateformes.

\textbf{Objectif.} Ce rapport documente les fondements stratégiques, organisationnels et réglementaires de \textbf{MABRID}, une place de marché numérique marocaine de confiance proposée. L'analyse exclut délibérément les détails d'implémentation et se concentre plutôt sur la vision commerciale, la gouvernance d'entreprise, l'architecture de sécurité, la stratégie cloud, la conformité juridique et la croissance durable.

\textbf{Méthodologie.} Le rapport synthétise des cadres de gestion établis (Business Model Canvas, Lean Startup, planification alignée sur le PMBOK), des normes internationales de sécurité (\gls{iso27001}, \gls{nistcsf}, contrôles \gls{cis}), des pratiques de gouvernance cloud sur Microsoft Azure, ainsi que le droit marocain de protection des données (loi 09-08) et les principes du \gls{gdpr} pertinents pour les traitements transfrontaliers.

\textbf{Résultats.} La différenciation concurrentielle de MABRID repose sur une confiance vérifiable, une gouvernance communautaire par ville, des politiques vendeurs transparentes et une posture de sécurité Zero Trust soutenue par une surveillance continue. Une feuille de route d'expansion nationale par phases équilibre la diversification des revenus (abonnements, services de vérification, partenariats d'entreprise) et la discipline des coûts dans les opérations cloud.

\textbf{Conclusions.} Le rapport démontre que MABRID peut être positionné comme une documentation prête pour l'investissement : un cas d'affaires cohérent, un programme de sécurité défendable et un modèle opérationnel orienté conformité, adapté à l'adoption par les entreprises et à l'examen réglementaire au Maroc.

\vspace{0.8cm}
\noindent\textbf{Mots-clés :} place de marché numérique, Maroc, gouvernance d'entreprise, Zero Trust, cloud Azure, loi 09-08, économie de plateforme, cybersécurité, conformité juridique, entrepreneuriat.
